\documentclass[11pt]{article}
\usepackage[margin=1in]{geometry}
\usepackage{amsmath,amssymb,amsthm,hyperref}
\newtheorem{lemma}{Lemma}\newtheorem{theorem}[lemma]{Theorem}
\newtheorem{corollary}[lemma]{Corollary}
\title{Ordinary gap-six terminal families with leading digit two}
\author{Internal proof; no priority or external-review claim}
\date{18 November 2026}
\begin{document}\setlength{\emergencystretch}{1em}\maketitle
This note treats the ordinary binomial $f(X)=X(A+BX^6)$,
where $A,B\in\mathbb C$ are not both zero.
At terminal order $N$, the assertion is that the two moments at
orders $N-1,N$ cannot both vanish.
For $AB\ne0$, their phase is the nonzero complex number$t=B/A$.
After removing nonzero scalar factors, these moments are
\[
 P_j(t)=\sum_{k=0}^j\binom jk(j+1)_{6k}t^k,\qquad j=N-1,N.
\]
If $f$ is a nonzero monomial its positive moments are nonzero,
so this case is immediate.
The theorem below proves all terminals $N=2\cdot7^r+c$ with
$0\le c\le6$, $r\ge0$. It does not prove all-index gap-six SFC,
all binomial SFC, or a Gamma-shape extension.

For the new cases fix $2\le c\le6$, $r\ge1$, $L=7^r$,
$N=2L+c$, and $C=N+5$.
Work in $\mathbb C_7$, with valuation ring$\mathcal O$,
maximal ideal$\mathfrak m$, and residue field$k$.
Let $\mathcal A=C-zD_z+z$,
$\Pi_N(X)=\prod_{\ell=0}^{N-1}(X+6\ell)$.
For $Q\in\mathcal O[X]$, $\deg Q\le4$, define
\[
 \begin{gathered}
 d_j(Q)=\sum_{i=0}^j(-1)^i\binom ji
       \binom{(C-i)/6+N-1}{N}Q(C-i),\\
 V_a(Q;v)=\sum_{\ell\ge0}\frac{7^\ell a!}{(a+6\ell)!}
                 d_{a+6\ell}(Q)v^\ell,\qquad0\le a<6.
 \end{gathered}                                                \tag{1}
\]
The finite formal identity
\[
 \frac{\Pi_N(\mathcal A)Q(\mathcal A)1}{6^NN!}
       =\sum_{j=0}^{N+4}d_j(Q)\frac{z^j}{j!}
\]
identifies these as normalized cyclic rows on$z^6=7v$.
Binomial polynomials are integer-valued on$\mathbb Z_7$ by
continuity from nonnegative integers, so$d_j(Q)\in\mathcal O$.
The cyclic coefficient weight has valuation
$(s_7(a+6\ell)-a)/6\ge0$.
Thus the rows are integral, and their basis-column matrix lies
in$\mathbb Z_7[v]$.

Put
\[
 S=S_r(v)=1+\sum_{i=1}^r(-1)^iv^{(7^i-1)/6},\quad
 h_c=\min(c+1,6),\quad
 \delta=v\partial_v+(S-1)\partial_S.
\]
For the physical seed set
$T_N(t)=\sum_j\binom Nj(N)_{6j}t^j$ and $H_N(Y)=7T_N(Y/7)$.
Rows$a<h_c$ will be divided once by7.
Write $W_a=\overline{V_a/7}$ there, and $W_a=\bar V_a$ for
$a\ge h_c$.

\begin{lemma}[Uniform formal seeds]\label{lem:seeds}
The indicated entire rows are seven-divisible on the five basis
columns. Their reduced seeds, in order$a=0,\ldots,5$, are
\begingroup\footnotesize\setlength{\arraycolsep}{3pt}
\[
\begin{array}{c|l}
c&(W_0(1),\ldots,W_5(1))\\\hline
2&(2,\ 2-S,\ 3S^2-3S+2,\ -3S^2-3S,\ -S^2,\ 0)\\
3&(2,\ S+2,\ -2S^2-2S+2,\ 2S^2-S+2,\ -S^2-3S,\ -3S^2)\\
4&(1-vS^2,\ 1-S,\ -2S^2+S+1,\ -3S^2+2S+1,\
     2S^2-3S+1,\ -3S^2+2S)\\
5&(3+vS-vS^2,\ 3+S+3vS^2,\
                 3+S+3S^2,\ 3+S+3S^2,\ 3+S+3S^2,\ 3+S+3S^2)\\
6&(0,\ 3vS^2+vS-S,\ S^2(1-v),\ 0,\ 0,\ 0).
\end{array}                                                    \tag{2}
\]
\endgroup
\end{lemma}
\begin{proof}
Use \texttt{lem:diagonal} and the formal calculation in the proof
of \texttt{lem:first-row} of
\href{GAP_SIX_TWO_POWER_PLUS_ONE_SFC.tex}{the two-power-plus-one note}.
The exact diagonal series is
\[
 \sum_{n\ge0}\frac{B_{n,n,6}(z)}{6^nn!}t^n
 =\kappa e^{z(1-R)},\quad R-R^7=t,\ R(0)=1,\quad
 \kappa=\frac{6R}{7R^6-1}\equiv R-7t\pmod{49}.
\]
For clarity we give the complete all-depth coefficient extraction.
Let $T=\sum_it^{7^i}$,
\[
 P(T)=T+3T^2+5T^3+5T^4+3T^5+T^6,\quad
 C(T)=\frac{T^7-\sum_it^{7^{i+1}}}{7}.
\]
Write $R=1+T+7R_1$, $R_1\equiv C(T)+P(T)\pmod7$.
Define integer polynomials$a_0=-1$,
$a_i=-1-va_{i-1}^7$, and$A_0=\sum_i a_it^{7^i}$.
Then$\bar a_r=-S_r(v)$.
The integral Artin--Hasse series from
\texttt{lem:artin-hasse-closure} of the
\href{ARITHMETIC_FRONTIER_TERMINAL_ORDERS.tex}{arithmetic note}
gives the integral formal parameter$A_*=A_0+7A_1$ with
\[
 e^{z(1-R)}=\mathcal E_7(zA_*),\quad z^6=7v,\qquad
 A_1\equiv-R_1-vC(A_0).
\]
The existence follows by recursive formal inversion of
$A_*+vA_*^7+7^6v^8A_*^{49}+\cdots=1-R$.
All operations are formal$t$-adic coefficient identities;
no exponential is evaluated at a7-adic boundary.

Set$m=c+2=s_7(N)$. For a divided seed$a<h_c$,
the leading section$(1+T)A_0^a$ uses at most$a+1<m$ selected
powers of7, so its coefficient of$t^N$ vanishes exactly.
The divided coefficient modulo7 is
\[
\begin{aligned}
 \relax[t^N]\bigl[&
 R_1A_0^{a-1}\{A_0-a(1+T)\}
 -av(1+T)A_0^{a-1}C(A_0)\\
 &+v\,a![X^{a+6}]\mathcal E_7(X)(1+T)A_0^{a+6}
                                  -tA_0^a\bigr]
\end{aligned}                                                \tag{3}
\]
interpreting the first term as$R_1$ at$a=0$ and omitting its
second term there. The cyclic expansion has only one additional
section modulo49. At$a=0$ its coefficient is$6v(1+T)A_0^6$.
For$a\ge1$, it contains$A_0^{a+6}=A_0^{a-1}A_0^7$.

There are no hidden intermediate-digit contributions.
For$m\le6$, every term in a divided seed has selected degree
at most$a+7\le m+5$; for$m=7,8$ it is at most12.
In both cases this is less than$m+6$, the minimum degree needed
for any carry representing$N$.
Thus every contributing term is carry-free and selects only
positions0 and$r$.
The term$tA_0^a$ has too few selected powers.
Set$U=x+y$, $B=-x-Sy$, and the integral cross polynomials
\[
 C_U=(U^7-x^7-y^7)/7,\qquad
 C_B=(B^7-(-x)^7-(-Sy)^7)/7 .
\]
The divided zero seed is
\[
 [x^cy^2]\{C_U+P(U)+6v(1+U)B^6\}.                            \tag{4}
\]
For$1\le a<h_c$ it is
\[
 [x^cy^2]\bigl[
 (C_U+P(U))B^{a-1}\{B-a(1+U)\}
                 -av(1+U)B^{a-1}C_B\bigr].                 \tag{5}
\]
Indeed the additional Artin--Hasse term at$a\ge1$ contains
$B^7=-x^7-S^7y^7$ in characteristic7 and cannot contribute
to$x^cy^2$, whose two exponents are below7.
The ordinary seeds$a\ge h_c$ are
$[x^cy^2](1+U)B^a$.
These fixed finite polynomial extractions give exactly (2).

Finally the exact raising recurrence is
\[
\begin{aligned}
 V_a(XQ)&=(C-a-6v\partial_v)V_a(Q)+aV_{a-1}(Q),\quad a\ge1,\\
 V_0(XQ)&=(C-6v\partial_v)V_0(Q)+7vV_5(Q)/5!.
\end{aligned}
\]
Starting with the divisible seeds, it propagates seven-divisibility
of all rows$a<h_c$ on$1,X,\ldots,X^4$.
\end{proof}

\begin{lemma}[The actual universal matrices]\label{lem:matrices}
The six-by-five matrix$\mathcal M_c(S,v)$ on$1,X,\ldots,X^4$
is generated from (2) by
\[
\begin{aligned}
 W_0(XQ)&=(c+5+\delta)W_0(Q)
                   +\begin{cases}vW_5(Q),&h_c<6,\\0,&h_c=6,\end{cases}\\
 W_a(XQ)&=(c+5-a+\delta)W_a(Q)
                   +\begin{cases}0,&a=h_c,\\aW_{a-1}(Q),&a\ne h_c.\end{cases}
                                                       \quad(1\le a<6).
\end{aligned}                                                 \tag{6}
\]
Its specialization is exactly the reduction of the actual matrix
after the stated global row divisions.
For$c=6$, row zero is identically zero; its only possibly nonzero
maximal minor is$m_0=S^5 f(S,v)$, obtained by omitting row zero.
\end{lemma}
\begin{proof}
Divide the indicated exact rows and reduce their raising recurrence.
Use$C\equiv c+5$, $-6\equiv1$, $5!\equiv1$, and$vS_r'=S_r-1$.
The first ordinary row has zero lower term, because its preceding
row is still seven-divisible.
When all six rows are divided, the wrap term disappears because
the ordinary reduction$\bar V_5$ is zero.
The zero row and determinant factor for$c=6$ follow by (2)--(6).
\end{proof}

\begin{lemma}[The actual moment residue]\label{lem:moment-residue}
Assume $c=6$. Set $U=S-1$ and
\[
 \mathcal Q(S,v)=1-v^7+2v(U+v)^6+v^2U^7(U+v)^5 .
\]
For every$r\ge1$, as polynomials in$v$ over$\mathbb F_7$,
\[
       v^N\overline H_N(-1/v)=\mathcal Q(S_r(v),v).                  \tag{7}
\]
\end{lemma}
\begin{proof}
The finite-field moment trace is \texttt{prop:trace} of the
\href{PHYSICAL_PRIME_TOTAL_DEGREE_REDUCTION.tex}
{physical prime-total-degree note}.
For $v\ne0$, the inverse eigenvalues$\tau$ at$Y=-1/v$ satisfy
\[
 P(\tau)=\tau^7+\tau/v-1=0,\qquad P'(\tau)=1/v.
\]
This polynomial is separable with nonzero roots.
Induction gives
\[
 \tau^{7^r}=b_r(\tau+S_r(v)-1),\qquad
 b_r=(-1)^rv^{-(7^r-1)/6}.
\]
Indeed $S_{r+1}=1-vS_r^7$.
Consequently
\[
 \overline H_N(-1/v)
      =b_r^{-2}\sum_{P(\tau)=0}\tau^{-6}(\tau+U)^{-2}.      \tag{8}
\]
Put $E=vS^7+U$.
At an actual$S=S_r(v)$,
\[
 E=S_r-S_{r+1}=(-1)^rv^{(7^{r+1}-1)/6}\ne0,
\]
so all denominators in (8) are nonzero.

Here is an explicit evaluation of the rational trace in (8).
Apply the residue theorem for rational functions to
\[
       t^{-6}(t+U)^{-2}\frac{P'(t)}{P(t)}.
\]
There is no residue at infinity. For$U\ne0$, the two other
residues, before multiplication by$1/v$, are
\[
 -\frac{(U+v)^5}{v^5U^7},\qquad
                 \frac{v^2S^7}{U^7E^2}.
\]
Negating their sum and multiplying by$1/v$ gives
\[
 \sum_{P(\tau)=0}\tau^{-6}(\tau+U)^{-2}
                         =\frac{\mathcal Q(S,v)}{v^6E^2}.          \tag{9}
\]
The cancellation uses the exact characteristic-seven identity
\[
 E^2(U+v)^5-v^7S^7=U^7\mathcal Q(S,v).
\]
After cancellation the only denominators are$v$ and$E$.
Thus (9) also specializes to$U=0$, where both are nonzero.
Alternatively it is a finite seven-dimensional algebra calculation:
$\tau^{-1}=\tau^6+1/v$,
$(\tau+U)^{-1}=
 [v\sum_{j=0}^6(-U)^j\tau^{6-j}+1]/E$,
and the traces of$1,\tau,\ldots,\tau^5$ are zero, while
$\operatorname{Tr}\tau^6=1/v$.
The checker independently evaluates this expression by reduction
modulo$P$.

Substitute the actual$E$ in (8)--(9).
Since $(7^{r+1}-1)/6=7(7^r-1)/6+1$,
the prefactor $b_r^{-2}/(v^6E^2)$ is exactly$v^{-N}$.
This proves (7) for$v\ne0$, and hence as a polynomial identity.
\end{proof}


\begin{lemma}[Complete necessary residue certificates]\label{lem:certificates}
The explicit arrays in
\href{gap_six_two_digit_leading_two_certificate.json}{the coefficient certificate}
have the following consequences over every residue extension field.
For$c=2,3,4$, a rank drop forces$q_c(S)=0$, $v=W_c(S)$.
The quotient recurrence$x_0=1$, $x_{j+1}=1-W_cx_j^7$ has complete
periods$114,1254,1092$ and every target$x_j-S$ is a unit.
The degrees of$q_c$ are$20,30,40$, respectively.

For$c=5$, every rank-drop point is either$(S,v)=(0,1)$ or a
point on$q_5(S)=0$, $v=W_5(S)$, $S\ne0$.
Here$\deg q_5=30$, the complete orbit has period72, and every
target is a unit.

For$c=6$, impose also the actual moment equation
$\mathcal Q(S,v)=0$ from Lemma~\ref{lem:moment-residue}.
Then a rank-drop point is either$(S,v)=(0,1)$ or a point on a
degree35 quotient with$v=W_6(S)$.
Its complete orbit has period728, and its only target gcds
are$S-1$, exactly at depths0 modulo7; also$W_6(1)=-1$.
\end{lemma}
\begin{proof}
The certificate retains each entire matrix, all six maximal minors,
every necessary relation with explicit sparse multiplier polynomials,
and every quotient-orbit state.
For$c=2,3,4$ the two relation targets are$q_c$ and$v-W_c$.
The quotient target gcd is1 at every state of the complete periods.

For$c=5$ the three necessary relation targets are
\[
 S^5q_5(S),\quad S^5v+h_5(S),\quad f_5(S,v),
                    \qquad f_5(0,v)=(v-1)^4.
\]
Thus$S=0$ forces$v=1$; otherwise divide by$S^5$.
The retained$W_5$ satisfies$S^5W_5+h_5=0$ modulo$q_5$,
and$S$ is a unit in that quotient.
The complete72-state record excludes$S=S_r(v)$ on this branch.

For$c=6$, at$S=0$ one has
$\mathcal Q(0,v)=-(v-1)^5$, so$v=1$.
At$S\ne0$, the determinant$m_0=S^5f$ forces$f=0$.
Explicit combinations of$f,\mathcal Q$ give
\[
 S^7q_6(S),\qquad S^7v+h_6(S).
\]
Here$S$ is a quotient unit and$S^7W_6+h_6=0$ modulo$q_6$.
The complete728-state certificate gives
\[
 \gcd(q_6,x_j-S)=
 \begin{cases}S-1,&j=0\pmod7,\\1,&\text{otherwise},\end{cases}
 \qquad W_6(1)=-1.
\]
For all cases the record includes the return to$x_0=1$,
which proves periodicity at every depth.
No field quotient, irreducibility, or inferred finite-orbit bound
is assumed.
The replay rebuilds each matrix from (4)--(6), recomputes every
minor by sparse permutation expansion, and verifies each
necessary combination by exact polynomial multiplication.
It does not trust a Groebner calculation.
\end{proof}

\begin{lemma}[The fixed residue packets]\label{lem:fixed-packets}
For$c=5,6$, the polynomial$H_N$ has no repeated root in the
residue disk$Y=-1$.
For$c=6$ and$r=0$ modulo7, $r\ge1$, it also has no repeated root
in the residue disk$Y=1$.
\end{lemma}
\begin{proof}
Use \texttt{prop:trace} of the
\href{PHYSICAL_PRIME_TOTAL_DEGREE_REDUCTION.tex}{physical-prime note}:
\[
 H_n(Y)\equiv K_n(720Y)\pmod{49},
\]
where$K_n(Z)$ is the power sum of the seven roots of
$(X-1)^7-ZX^6$.
For$Y=-1$, $Z=-720\equiv15\pmod{49}$.
By \texttt{lem:trace-parity} of the
\href{GAP_SIX_SMALL_OFFSET_FAMILIES.tex}{small-offset note},
these point values have period two in$r\ge1$.
Newton identities give the exact representatives
\[
\begin{array}{c|cc}
c&H_N(-1),\ r\text{ odd}&H_N(-1),\ r\text{ even}\\\hline
5&21&29\\
6&7&37
\end{array}\qquad\pmod{49}.                                \tag{10}
\]
They are checked at$N=19,103,20,104$.
For even$r$ the disk has no root; for odd$r$ the point value
has valuation one.
By \texttt{lem:local-cluster} of the same small-offset note,
the full Hensel factor in that disk is Eisenstein and hence
separable. This works for every local width and every ramified
point, without any odd-degree assumption.

For$Y=1$, work in the finite free algebra with
\[
 P_{34}(X)=(X-1)^7-34X^6
 =X^7+8X^6+21X^5+14X^4+35X^3+28X^2+7X+48\pmod{49}.
\]
Its inverse class$\tau$ satisfies$\tau^7-\tau=1$ modulo7.
Seven Frobenius iterations give$\tau^{7^7}=\tau$ and hence
$X^{7^7}=X$ modulo7.
Taking seventh powers of a lift gives$X^{7^8}=X^7$ modulo49.
Raising to$7^{r-1}$ proves period seven in$r\ge1$ for
$K_{2\cdot7^r+6}(34)$.
At the reference$r=7$, the remainder of$X^{2\cdot7^7+6}$ is
$[20,29,21,7,28,7,15]$ in ascending degree.
Newton traces of$1,X,\ldots,X^6$ are
$[7,41,22,48,8,34,15]$.
Their dot product is21 modulo49.
Thus for every positive$r$ divisible by7,
$H_N(1)\equiv21\pmod{49}$, again giving an Eisenstein and
separable full local factor.
The coefficient certificate and replay verify both polynomial
period identities and all the stated point traces.
\end{proof}

\begin{theorem}[The leading-two terminal family]\label{thm:family}
Every nonzero ordinary binomial $X(A+BX^6)$ satisfies its consecutive
moment test at every terminal
\[
                    N=2\cdot7^r+c,\qquad0\le c\le6,\quad r\ge0.
\]
\end{theorem}
\begin{proof}
For$c=0$ use \texttt{cor:pair} of the physical-prime note:
the terminal index is in its native Eisenstein family.
For$c=1,r\ge1$ use \texttt{thm:family} of
\href{GAP_SIX_TWO_POWER_PLUS_ONE_SFC.tex}{the two-power-plus-one note}.
At$r=0$, $N=2,\ldots,7$ are native; $N=8$ is covered by
\texttt{cor:scaled-family}, with$b=1,r=1$, of
\href{GAP_SIX_SCALED_POWER_PLUS_ONE_SFC.tex}{the scaled-family note}.
This last endpoint is not claimed to be native.

It remains to treat$2\le c\le6,r\ge1$.
Put$T_N(t)=\sum_j\binom Nj(N)_{6j}t^j$.
The exact lower-one bridge gives
\[
 P_{N-1}=T_N-\theta T_N/N,\qquad
 P_N=T_N+6\theta T_N/N.
\]
The determinant is$7/N\ne0$ and$T_N(0)=1$; thus a common phase
would be a multiple root of$T_N$.
By \texttt{thm:radius} of the physical-prime note it has
$u=-1/t=7v$ with$v$ a unit.

The necessary cyclic rank relation is
\texttt{lem:critical-ideal} of the
\href{gap_three_hurwitz_structure.tex}{adjoint-structure note}.
Its complex statement transfers at each algebraic multiple root:
the cyclic minors for fixed rational$N,\gamma=N$ are rational
polynomials, so apply the complex theorem and transport their
vanishing through algebraic embeddings.
No exponential is evaluated on the7-adic boundary.
At the fixed shape,
$B_{N+j}=\Pi_N(\mathcal A)
\prod_{\ell=N}^{N+j-1}(\mathcal A+6\ell)1$.
These five monic multiplier polynomials have consecutive degrees,
so the columns differ from$Q=1,X,\ldots,X^4$ by a unit
triangular change.
Normalize a relation vector to primitive integral coefficients,
divide the indicated entire rows, and reduce.
This gives a nonzero kernel of$\mathcal M_c(S_r(\bar v),\bar v)$.

For$c=2,3,4$, the complete unit-target cycles exclude it.
For$c=5$ the only actual residue point is$v=1$; hence
$Y=7t=-1/v$ lies in the disk$Y=-1$, excluded by
Lemma~\ref{lem:fixed-packets}.
For$c=6$, the same actual root also satisfies$H_N(-1/v)=0$,
so Lemma~\ref{lem:moment-residue} imposes$\mathcal Q=0$.
The complete coupled certificate leaves either$v=1$, or
$v=-1$ with$r=0$ modulo7.
The former enters$Y=-1$; the latter enters$Y=1$ at exactly
the depths covered by Lemma~\ref{lem:fixed-packets}.
Neither disk contains a repeated root.
This contradiction proves the new cases.
\end{proof}

\begin{corollary}[The refined degree-seven frontier]
\label{cor:degree-seven-frontier}
Put
\[
\begin{aligned}
 \mathcal U_2={}&\{b7^s:1\le b\le6,\ s\ge0\}
 \cup\{b7^r+1:1\le b\le6,\ r\ge0\}\\
 &{}\cup\{7^r+c:2\le c\le6,\ r\ge0\}
 \cup\{2\cdot7^r+c:2\le c\le6,\ r\ge0\}.
\end{aligned}
\]
Ordinary binomial SFC through degree seven is equivalent to
$\gcd_{\mathbb Q[t]}(P_{N-1},P_N)=1$ for every terminal$N\ge6$
outside$\mathcal U_2$.
Any hypothetical common phase on those remaining tests has
$v_7(-1/t)=1$ at every algebraic embedding.
\end{corollary}
\begin{proof}
Use \texttt{cor:degree-seven-frontier} of the small-offset note
and remove the additional terminals supplied by
Theorem~\ref{thm:family}.
The converse is restriction to a subset of the original tests.
The necessary radius is retained from the prior frontier.
Infinitely many terminal orders remain; neither degree-seven
SFC nor all-binomial SFC is closed by this corollary.
\end{proof}

The exact replay verifies polynomial identities and complete
finite certificates. It is not formal verification or external review.
No Gamma-shape extension or statement for arbitrary auxiliary seeds
is inferred.
\end{document}
